\subsection{单模的张量积}

设 \(A_{1}\) 与 \(A_{2}\) 是交换域 K 上的代数。我们把 K-代数 \(A_{1} \otimes A_{2}\) 记作 A。

引理 1. —— 设 \(M_{1}\) 与 \(N_{1}\) 是 \(A_{1}\)-模，\(M_{2}\) 与 \(N_{2}\) 是 \(A_{2}\)-模。我们作如下假设：

\begin{enumerate}
\def\labelenumi{(\roman{enumi})}
\item
  \(A_{1}\)-模 \(M_{1}\) 是有限生成的。
\item
  \(A_{2}\)-模 \(M_{2}\) 是有限生成的，或者 \(N_{1}\) 在 K 上有限维。
\end{enumerate}

令 \(\mathrm{M} = \mathrm{M}_1\otimes \mathrm{M}_2\) 与 \(\mathrm{N} = \mathrm{N}_1\otimes \mathrm{N}_2\)，并把它们视为环 \(\mathrm{A} = \mathrm{A}_1\otimes \mathrm{A}_2\) 上的模。典范同态（II, §3, No.~5, 第 251 页）

\[
\lambda : \operatorname{Hom} _ {\mathrm{K}} (\mathrm{M} _ {1}, \mathrm{N} _ {1}) \otimes \operatorname{Hom} _ {\mathrm{K}} (\mathrm{M} _ {2}, \mathrm{N} _ {2}) \longrightarrow \operatorname{Hom} _ {\mathrm{K}} (\mathrm{M}, \mathrm{N})
\]

于是诱导 K-向量空间的一个同构

\[
\varphi : \operatorname{Hom} _ {\mathrm{A} _ {1}} (\mathrm{M} _ {1}, \mathrm{N} _ {1}) \otimes \operatorname{Hom} _ {\mathrm{A} _ {2}} (\mathrm{M} _ {2}, \mathrm{N} _ {2}) \longrightarrow \operatorname{Hom} _ {\mathrm{A}} (\mathrm{M}, \mathrm{N}).
\]

映射 \(\lambda\) 是单射（II, §7, No.~7, 第 308 页，命题 16），并把线性子空间 \(\mathrm{Hom}_{\mathrm{A}_1}(\mathrm{M}_1,\mathrm{N}_1)\otimes \mathrm{Hom}_{\mathrm{A}_2}(\mathrm{M}_2,\mathrm{N}_2)\) 映到 \(\mathrm{Hom}_{\mathrm{A}}(\mathrm{M},\mathrm{N})\)。因此只需证明每个从 M 到 N 的 A-线性映射都属于 \(\mathrm{Hom}_{\mathrm{A}_1}(\mathrm{M}_1,\mathrm{N}_1)\otimes \mathrm{Hom}_{\mathrm{A}_2}(\mathrm{M}_2,\mathrm{N}_2)\) 在 \(\lambda\) 下的像。设 \(u:\mathrm{M}\to \mathrm{N}\) 是一个 A-线性映射。设 \(x\in \mathrm{M}_1\)。用 \(u_{x}\) 表示 \(\mathrm{A}_2\)-线性映射 \(y\mapsto u(x\otimes y)\)（从 \(\mathrm{M}_2\) 到 \(\mathrm{N}_1\otimes \mathrm{N}_2\)）。令 \(\mathrm{P} = \mathrm{Hom}_{\mathrm{A}_2}(\mathrm{M}_2,\mathrm{N}_2)\)，并用 \(\nu\) 表示从 \(\mathrm{N}_1\otimes \mathrm{P}\) 到 \(\mathrm{Hom}_{\mathrm{A}_2}(\mathrm{M}_2,\mathrm{N}_1\otimes \mathrm{N}_2)\) 的典范同态（II, §4, No.~2, 第 269 页）。这个映射是单射（II, §4, No.~2, 第 269 页，命题 2, (i)，应用于 K-向量空间 \(\mathbf{N}_1\)）。由假设 (ii)，存在一个线性子空间 \(\mathrm{V}_x\)（\(\mathbf{N}_1\) 的，在 K 上有限维），使得 \(u_{x}\) 取值于 \(\mathrm{V}_x\otimes \mathrm{N}_2\)。由此可知，\(u_{x}\) 是 \(\nu\) 下某个唯一元素 \(v_{x}\)（属于 \(\mathbf{N}_1\otimes \mathbf{P}\)）的像。映射 \(\widetilde{u}:x\mapsto v_x\)（从 \(\mathbf{M}_1\) 到 \(\mathbf{N}_1\otimes \mathbf{P}\)）是 \(\mathbf{A}_1\)-线性的。由假设 (i)，\(\mathbf{A}_1\)-模 \(\mathbf{M}_1\) 是有限生成的。与上面类似的论证表明，\(\widetilde{u}\) 属于 \(\mathrm{Hom}_{\mathrm{A}_1}(\mathrm{M}_1,\mathrm{N}_1)\otimes \mathrm{P}\) 在 \(\mathrm{Hom}_{\mathrm{A}_1}(\mathrm{M}_1,\mathrm{N}_1\otimes \mathrm{P})\) 中的像。引理 1 得证。

定理 2. —— 设 \(A_{1}\) 与 \(A_{2}\) 是交换域 K 上的代数；设 \(S_{1}\) 是单 \(A_{1}\)-模，\(S_{2}\) 是单 \(A_{2}\)-模。设 \(D_{1}\) 与 \(D_{2}\) 分别是 \(S_{1}\) 与 \(S_{2}\) 的交换子。令 \(M = S_{1} \otimes S_{2}\)，\(A = A_{1} \otimes A_{2}\)，\(D = D_{1} \otimes D_{2}\)。我们把 M 视为左 (A,D)-双模。

\begin{enumerate}
\def\labelenumi{\alph{enumi})}
\item
  A-模 M 的交换子等于 \(\mathrm{D}_{\mathrm{M}}\)。
\item
  映射 \(a \mapsto aM\) 是从按包含关系排序的 D 的右理想所成之集到按包含关系排序的 M 的 A-子模所成之集的同构。逆映射把 M 的子模 N 映到由满足 \(dM \subset N\) 的元素 d 组成的 D 的理想。
\end{enumerate}

断言 a) 由引理 1 得出，因为单模是单生成的。

设 \(\mathrm{T}\) 是 \((\mathrm{A}_1, \mathrm{D}_2)\)-双模 \(\mathrm{S}_1 \otimes (\mathrm{D}_2)_d\)。我们把 \(\mathrm{M} = \mathrm{S}_1 \otimes \mathrm{S}_2\) 与 \(\mathrm{T} \otimes_{\mathrm{D}_2} \mathrm{S}_2\) 视为同一（II, §3, No.~8, 第 258 页）；这一视为同一与环 \(\mathrm{A} = \mathrm{A}_1 \otimes \mathrm{A}_2\) 上的左模结构相容。

设 \(\mathrm{N}\) 是 \(\mathrm{M}\) 的一个 A-子模；它是 \(\mathrm{A}_2\)-子模（属于 \(\mathrm{T} \otimes_{\mathrm{D}_2} \mathrm{S}_2\)），在形如 \((a_1)_\mathrm{T} \otimes 1_{\mathrm{S}_2}\)（其中 \(a_1\) 取遍 \(\mathrm{A}_1\)）的自同态下稳定。由 VIII, 第 63 页，推论 2，存在唯一的 \((\mathrm{A}_1, \mathrm{D}_2)\)-子双模 \(\mathrm{V}\)（\(\mathrm{T}\) 的），使得 \(\mathrm{N} = \mathrm{V} \otimes_{\mathrm{D}_2} \mathrm{S}_2\)。

同构 \(u\)（从 \(\mathrm{T} = \mathrm{S}_1 \otimes (\mathrm{D}_2)_d\) 到 \(((\mathrm{D}_1)_d \otimes (\mathrm{D}_2)_d) \otimes_{\mathrm{D}_1} \mathrm{S}_1\)，由 \(u(s \otimes d) = 1 \otimes d \otimes s\) 定义）是 \((\mathrm{A}_1, \mathrm{D}_2)\)-线性的。我们把这些 \((\mathrm{A}_1, \mathrm{D}_2)\)-双模视为同一。与上面给出的论证类似，可以证明存在唯一的右 \((\mathrm{D}_1 \otimes \mathrm{D}_2)\)-子模 \(\mathfrak{a}\)（\(\mathrm{D}_1 \otimes \mathrm{D}_2\) 的），使得 \(\mathrm{V} = \mathfrak{a} \otimes_{\mathrm{D}_1} \mathrm{S}_1\)。考虑到上面所作的视为同一，\(\mathfrak{a}\) 是 \(\mathrm{D} = \mathrm{D}_1 \otimes \mathrm{D}_2\) 的唯一使得 \(\mathrm{N} = \mathfrak{a}\mathrm{M}\) 的右理想。

我们刚刚证明了映射 \(a \mapsto aM\) 是双射；最后的断言由此得出。

推论 1. —— 模 \(S_{1} \otimes S_{2}\)（环 \(A_{1} \otimes A_{2}\) 上的）是半单的（相应地，同型的、单的）当且仅当环 \(D = D_{1} \otimes D_{2}\) 是半单的（相应地，单环、体）。特别地，\(S_{1} \otimes S_{2}\) 是单模（当 \(S_{1}\) 或 \(S_{2}\) 的交换子等于 K 时）。

由定理 2，模 \(S_{1} \otimes S_{2}\)（环 \(D_{1} \otimes D_{2}\) 上的）是半单的（相应地，同型的、单的）当且仅当右 D-模 \((\mathrm{D}_{1} \otimes \mathrm{D}_{2})_{d}\) 是（VIII, 第 109 页，命题 10）。而右 D-模 \(D_{d}\) 是单模当且仅当 D 是体；它是同型的（相应地，半单的）当且仅当环 D 是单环（相应地，半单环）（VIII, 第 120 页，定义 1；VIII, 第 121 页，推论 1；以及 VIII, 第 137 页，命题 2）。

推论 2. —— 我们有 \(\mathfrak{R}_{\mathrm{A}}(\mathrm{M}) = \mathfrak{R}(\mathrm{D})\mathrm{M}\)。A-模 M 无根当且仅当环 D 无根。

这由 VIII, 第 108 页，命题 8 与定理 2, b) 得出。

